\section{Study case failure signature}

This chapter focuses on the functionnal failure caused by an \gls{esd} of an integrated function.
The studied function is the primary supply of a complex \gls{asic}.
This supply plays a critical role in the functionning of the entire product.
It is connected to the battery of the vehicle, and is the first block of the product to start.
It wakes up and powers all other functions inside the integrated circuit.
In case of failure, it can cause the entire product to shut down.

SCHEMATIC OF THE TESTCHIP

The primary supply processes the battery supply in a chained fashion.
A first block (pre-regulator) clamps the battery voltage (that can reach up to 40V) to 9V, a more acceptable voltage for the silicium technology.
This clamped voltaged is then used to power up a bandgap reference.
Once properly started, this bandgap generates a 1.23 V voltage reference, stable accross a wide range of temperature, process variation and mismatchs.
The bandgap also outputs a 10uA current reference, stable in the same conditions.

After the bandgap, a \gls{ldo} regulator relies on the voltage reference to generate a stable 2.5V supply voltage, able to deliver and sustain up to 20mA.
This first regulator is connected externally to a 100nF decoupling capacitor to absorb peak currents and achieve stability.
This supply is used to supply most analog (digital ?) functions inside the integrated circuit.

A second regulator performs the same task, but starts with a delay compared to the first one.
It also ouputs a 2.5V supply, that powers digital logic inside the circuit.

\subsection{Full function failure modeling with a bottom-up approach}

Consider a system constituted of two blocks A and B connected together.

\begin{figure}[h]
  \centering
  \includegraphics{src/4/figures/system_ab.pdf}
  \caption{Basic IC function}
  \label{basic_ic_function}
\end{figure}

During normal operation, the \textbf{\textit{out}} signal complies with its specification,
under the condition that the \textbf{\textit{in}} signal does too.
This guarantees that block A and B are performing as expected.
When the system is exposed to an ESD on the input during normal operation, the \textbf{\textit{in}} signal is disturbed.
It no longer complies with its specification for a short amount of time.

The robustness of the entire system is defined by its ability to maintain signal \textbf{\textit{out}} in specification
while signal \textbf{\textit{in}} is temporarily disturbed.
Time is a key parameter here. On one extreme side, a disturbance that lasts forever is equivalent to the \textbf{\textit{in}} signal not fullfilling its DC specification,
and thus the \textbf{\textit{out}} signal will not fullfill its spec as well.
On the other extreme side, an extremely short disturbance, much shorter than the functionnal timescales of the integrated function will most likely not disturb ??

The first task for defining the robustness of an integrated function is to determine the time threshold under which it can maintain functionnality while being disturbed.
The second key parameter is related to the amplitude of the disturbance, whether it is a voltage, current, energy, etc.
A larger amplitude can make a short pulse more harmful than a lomg one.
